\begin{prob}
    Consider the version of binary search shown below. This is the same as
    the code given in lecture, except that a single \python{print} statement
    has been added. This \python{print} statement prints the value of
    \python{stop} at the beginning of each call to \python{binary_search}.

    \begin{minted}[autogobble]{python}
        import math
        def binary_search(arr, t, start, stop):
            """
            Searches arr[start:stop] for t.
            Assumes arr is sorted.
            """
            print(stop) # <-- added print statement
            if stop - start <= 0:
                return None
            middle = math.floor((start + stop)/2)
            if arr[middle] == t:
                return middle
            elif arr[middle] > t:
                return binary_search(arr, t, start, middle)
            else:
                return binary_search(arr, t, middle+1, stop)
    \end{minted}

    Suppose \python{binary_search} is called on a \textbf{sorted} list with
    \python{start = 0} and \python{stop = len(n)}. True or False: the numbers
    printed must be a non-increasing sequence. That is, if $x$ is printed
    before $y$, then it must be the case that $x \geq y$.

    \Tf{}

\end{prob}
